\documentclass[11pt]{article}
\usepackage[T1]{fontenc}
\usepackage{lmodern}
\usepackage{amsmath,amssymb,amsthm,mathtools}
\usepackage[letterpaper,margin=1in]{geometry}
\usepackage{microtype}
\usepackage{enumitem}
\usepackage{booktabs}
\usepackage{graphicx}
\usepackage{tikz}
\usetikzlibrary{arrows.meta,calc,positioning}
\usepackage{xcolor}
\definecolor{linkblue}{HTML}{23527C}
\definecolor{parentblue}{HTML}{28698D}
\definecolor{childorange}{HTML}{B76223}
\usepackage[colorlinks=true,linkcolor=linkblue,citecolor=linkblue,urlcolor=linkblue]{hyperref}
\usepackage[capitalise,noabbrev]{cleveref}
\hypersetup{pdftitle={Thompson's group F is nonamenable}}

\newcommand{\R}{\mathbb R}
\newcommand{\N}{\mathbb N}
\newcommand{\Z}{\mathbb Z}
\newcommand{\norm}[1]{\lVert #1\rVert}
\newcommand{\ip}[2]{\langle #1,#2\rangle}
\newcommand{\id}{\mathrm{id}}
\DeclareMathOperator{\Lip}{Lip}
\DeclareMathOperator{\dist}{dist}
\DeclareMathOperator{\sinc}{sinc}
\newcommand{\unverified}[1]{\par\smallskip\noindent\fbox{\parbox{0.95\linewidth}{\textbf{Unverified step.} #1}}\par\smallskip}
\newtheorem{theorem}{Theorem}[section]
\newtheorem{lemma}[theorem]{Lemma}
\newtheorem{proposition}[theorem]{Proposition}
\newtheorem{corollary}[theorem]{Corollary}
\theoremstyle{definition}
\newtheorem{definition}[theorem]{Definition}
\newtheorem{remark}[theorem]{Remark}
\numberwithin{equation}{section}
\setlist[itemize]{topsep=4pt,itemsep=2pt,parsep=0pt}
\setlist[enumerate]{topsep=4pt,itemsep=2pt,parsep=0pt}
\setlength{\emergencystretch}{2em}

\title{Thompson's group \(F\) is nonamenable}
\author{Rewritten for readability from the OpenAI manuscript of the same title}
\date{}

\begin{document}
\maketitle
\begin{abstract}
We prove that Thompson's group \(F\) is nonamenable. This confirms
Geoghegan's conjecture and resolves the amenability problem for \(F\).
\end{abstract}

\section{The result and the plan of proof}\label{sec:introduction}

\subsection{Statement}

Thompson's group $F$ consists of the increasing piecewise linear
homeomorphisms of $[0,1]$ with finitely many pieces, dyadic rational
breakpoints, and slopes in $2^{\Z}$. Products are compositions,
$hg=h\circ g$, and $F$ is regarded as a discrete group.

A discrete group $\Lambda$ is \emph{amenable} if there is a linear
functional $\mu:\ell^\infty(\Lambda;\R)\to\R$ with $\mu(1)=1$,
$\mu(\varphi)\ge0$ whenever $\varphi\ge0$, and
$\mu(g\mapsto\varphi(hg))=\mu(\varphi)$ for every $h\in\Lambda$
\cite{vonNeumann1929,Day1957}. We use only the forward implication of
the F\o lner criterion \cite{Folner1955}: if $\Lambda$ is amenable, then
for every finite $S\subset\Lambda$ and every $\varepsilon>0$ there is a
nonempty finite $A\subset\Lambda$ with
\begin{equation}\label{eq:folner}
  \frac{|hA\mathbin{\triangle}A|}{|A|}<\varepsilon\qquad(h\in S).
\end{equation}
(The left-translation formulation in \cite[Theorem~1.2]{Moore2013}
bounds the sum of these ratios and hence implies each strict individual
bound by taking a smaller tolerance.)

\begin{theorem}\label{thm:main}
Thompson's group \(F\) is not amenable.
\end{theorem}

The only other outside input is the following.

\begin{lemma}[Benyamini--Sternfeld]\label{lem:displacement}
There are a real Hilbert space \(H\), a Lipschitz map \(f:B\to B\) on its
closed unit ball, and a constant \(\delta>0\) such that
\(\norm{f(x)-x}\ge\delta\) for all \(x\in B\).
\end{lemma}

Benyamini and Sternfeld prove this for every infinite-dimensional normed
space \cite[Theorem, part~(3), p.~439]{BenyaminiSternfeld1983}.
Appendix~\ref{app:auxiliary} gives a self-contained construction on
$L^2([0,1];\R^2)$ with $\delta=1/2$. In finite dimensions the lemma is
false (Brouwer), so the infinite-dimensional space is essential.

\paragraph{Context.}
Thompson introduced $F$ in 1965; an early published construction is in
\cite{McKenzieThompson1973}, with the identification explained in
\cite[pp.~215--216]{CannonFloydParry1996}. Geoghegan's 1979 conjecture
that $F$ is nonamenable is recorded in
\cite[p.~227]{CannonFloydParry1996}. The problem is delicate because $F$
has no nonabelian free subgroup \cite[Theorem~3.1]{BrinSquier1985} yet is
not elementary amenable \cite[Theorem~4.10]{CannonFloydParry1996}; it
also has a classifying space with finitely many cells in each dimension
\cite{BrownGeoghegan1984}. Earlier constraints on amenability come from
sizes of F\o lner sets \cite[Theorem~1.1]{MooreFolner2013}, a convex
Ramsey characterization for binary trees \cite[Theorem~3.1]{Moore2013},
densities and diagram classes \cite{Guba2025Density,Guba2026}, and the
reduced $C^*$-algebra of Thompson's group $T$
\cite[Theorem~5.5]{HaagerupOlesen2017}. Earlier claimed resolutions:
nonamenability \cite{Akhmedov2021}; amenability \cite{Shavgulidze2009},
with errors identified in \cite{Moore2011Critique}; and a withdrawn
amenability claim \cite{Moore2012Withdrawn} (a different work from
\cite{Moore2013}). None of this is used below.

\subsection{Proof at a glance}

Fix $f$, a Lipschitz constant $L$ and $\delta$ from
Lemma~\ref{lem:displacement}, and an integer $D>4L^2/\delta^2$.

\emph{Objects.} A basic interval is $[k2^{-r},(k+1)2^{-r}]$; a basic
partition is a finite partition of $[0,1]$ into basic intervals. Fix
$D$ basic intervals $I_1<\dots<I_D$ (``parents'') with gaps between
them and endpoints in $(0,1)$, and inside each parent $I_i$ the affine
copies $I_i\cdot I_j$ of all parents (``children''). Every basic
partition $T$ gets a \emph{color} $p(T)\in B$ by recursion on the number
of cells: $p(T)$ is $f$ applied to the mean of the colors of the $D$
partitions obtained by restricting $T$ to a parent and rescaling to
$[0,1]$. For $g\in F$ and $n$ large, the image $gT^{(n)}$ of the uniform
partition of mesh $2^{-n}$ is basic, and $X_I(g)\in B$ denotes the color
of its rescaled restriction to $I$.

\emph{Key idea.} These vectors have two exact properties:
\begin{enumerate}[label=(\alph*)]
\item \emph{Recursion:} $X_{I_i}(g)=f(z_i(g))$, where
$z_i(g)=\frac1D\sum_jX_{I_i\cdot I_j}(g)$.
\item \emph{Transport:} for intervals $I<J$ of the family some
$h_{I,J}\in F$ maps $I,J$ affinely onto $I_1,I_2$, and then
$X_I(g)=X_{I_1}(h_{I,J}g)$, $X_J(g)=X_{I_2}(h_{I,J}g)$.
\end{enumerate}
Let $S$ be the finite set of all $h_{I,J}$; it is fixed before any
finite set $A\subset F$ is considered. Put
$\eta=\max_{h\in S}|hA\triangle A|/|A|$. By (b), the $A$-average of
$\ip{X_I}{X_J}$ is within $\eta$ of one number $\alpha$ for \emph{all}
separated pairs. Expanding the $A$-average of $\norm{z_i-m}^2$, where
$m=\frac1D\sum_kX_{I_k}$, the $\alpha$'s cancel and the non-separated
pairs form a fraction $1/D$ of the terms, so it is at most $4/D+4\eta$.
By (a), $m$ is the mean of the $f(z_i)$, so $\norm{m-f(m)}^2\le
\frac{L^2}D\sum_i\norm{z_i-m}^2$. Averaging,
$\delta^2\le L^2(4/D+4\eta)$: a positive lower bound for $\eta$
independent of $A$, contradicting \eqref{eq:folner}.

\emph{Dependencies.}
\begin{itemize}
\item \S\ref{sec:dyadic}: restriction composes \eqref{dyad:restriction};
fine images are basic (Lemma~\ref{dyad:eventual}); pair transport
(Lemma~\ref{dyad:transport}); covariance \eqref{dyad:covariance}.
\item \S\ref{sec:coloring}: the coloring, using $f$.
\item \S\ref{sec:boundary}: Proposition~\ref{avg:boundary} (the bound on
$\eta$). Property (b) uses Lemmas~\ref{dyad:eventual},
\ref{dyad:transport} and \eqref{dyad:covariance}; (a) uses
\eqref{dyad:restriction}; the contradiction uses
Lemma~\ref{lem:displacement}.
Theorem~\ref{thm:main} follows from Proposition~\ref{avg:boundary} and
\eqref{eq:folner}.
\item Appendix~\ref{app:auxiliary} proves Lemma~\ref{lem:displacement};
\S\ref{sec:consequences} (consequences) is not used in the proof.
\end{itemize}

\emph{Renamed symbols.} The source reuses several letters. Here the
generic group is $\Lambda$ and its mean $\mu$; in
\S\ref{sec:consequences} the generating set is $\Sigma$, the Hilbert
space $\mathcal K$, the conjugating operator $V$, the connection kernel
$\mathcal T_p$; in the appendix $\Lip(\gamma')$ is $K_\gamma$ and the
rotated field $Jv$ is written $v^\perp$.

\section{Recursive colors and finite averages}\label{sec:proof}

\subsection{Dyadic partitions and affine transport}\label{sec:dyadic}

We use the standard dyadic-partition model of $F$
\cite[Section~2, Lemma~2.2]{CannonFloydParry1996}.

A \emph{basic interval} is $[k2^{-r},(k+1)2^{-r}]$ with integers
$r\ge0$, $0\le k<2^r$. A \emph{basic partition} is a finite set of basic
intervals (its \emph{cells}) covering $[0,1]$ with pairwise disjoint
interiors; its \emph{mesh} is the largest cell length. A partition $T$
\emph{respects} an interval $I$ if $I$ is a union of cells of $T$. We
use the standard fact that two basic intervals whose interiors meet are
nested (dyadic intervals form a binary tree under inclusion).

For a basic interval $I=[a2^{-r},(a+1)2^{-r}]$ let
$s_I(x)=2^{-r}(x+a)$, the increasing affine bijection $[0,1]\to I$, and
write $I\cdot J=s_I(J)$. If $T$ respects $I$, its \emph{normalized
restriction} is
\[
  T_I=\{s_I^{-1}(C):C\in T,\ C\subseteq I\}.
\]

\emph{Basicness.} $s_I$ maps $[b2^{-t},(b+1)2^{-t}]$ onto
$[(a2^t+b)2^{-(r+t)},(a2^t+b+1)2^{-(r+t)}]$, so $I\cdot J$ is basic. A
basic $C\subseteq I$ is $[k2^{-s},(k+1)2^{-s}]$ with $s\ge r$, and
$s_I^{-1}(C)$ has left endpoint $(k-a2^{s-r})2^{-(s-r)}$ and length
$2^{-(s-r)}$, so it is basic. Hence $T_I$ is a basic partition.

\emph{Restriction composes.} Both $s_{I\cdot J}$ and $s_I\circ s_J$ are
the increasing affine bijection $[0,1]\to I\cdot J$, hence equal.
Suppose $T$ respects $I$. Applying $s_I$ shows that $J$ is a union of
cells $s_I^{-1}(C)$ of $T_I$ exactly when $I\cdot J$ is a union of cells
$C$ of $T$. Thus $T_I$ respects $J$ if and only if $T$ respects
$I\cdot J$, and then
\begin{equation}\label{dyad:restriction}
  (T_I)_J=\{s_J^{-1}s_I^{-1}(C):C\in T,\ C\subseteq I\cdot J\}=T_{I\cdot J}.
\end{equation}

Let $T^{(n)}$ be the uniform partition into intervals of length
$2^{-n}$, and for $g\in F$ let $gT^{(n)}$ be the set of images of its
cells.

\begin{lemma}\label{dyad:eventual}
For each $g\in F$, the partitions $gT^{(n)}$ are basic for all
sufficiently large $n$, and their meshes tend to zero. In particular,
given a finite set $K\subset F$ and a finite family $\mathcal J$ of
basic intervals, there is an $n_0$ such that $gT^{(n)}$ is basic and
respects every member of $\mathcal J$ whenever $g\in K$ and $n\geq n_0$.
\end{lemma}

This lemma lets one level $n$ serve all finitely many group elements in
Proposition~\ref{avg:boundary}.

\begin{proof}
\emph{Step 1.} Every linear piece of $g$ is $x\mapsto2^qx+b$ with
$q\in\Z$ and $b$ dyadic: starting from $g(0)=0$, the image of each
breakpoint is a sum of increments (power of two times dyadic length),
hence dyadic.

\emph{Step 2.} Write $b=c2^{-N}$ with $c\in\Z$, $N\ge0$. Let $n$ be so
large that all breakpoints are multiples of $2^{-n}$ and $n\ge q+N$ on
every piece. Each cell $[k2^{-n},(k+1)2^{-n}]$ then lies in one piece
and its image is
$[(k+c2^{n-q-N})2^{q-n},(k+c2^{n-q-N}+1)2^{q-n}]\subseteq[0,1]$, a basic
interval of length $2^{q-n}$. So $gT^{(n)}$ is basic with mesh at most
$2^{q_{\max}-n}\to0$.

\emph{Step 3.} A basic partition with mesh smaller than the length of a
basic interval $I$ respects $I$: every cell whose interior meets that of
$I$ is nested with $I$ and, being shorter, lies in $I$; these cells
cover $I$.

\emph{Step 4.} Apply Step 3 to each member of $\mathcal J$ and take the
largest of the finitely many thresholds over $g\in K$.
\end{proof}

Write $I<J$ when $\max I<\min J$ (a positive gap is required). An
interval is \emph{internal} if both endpoints lie in $(0,1)$.

\begin{lemma}\label{dyad:transport}
If $I<J$ and $I'<J'$ are pairs of internal basic intervals, there is
an $h\in F$ whose restrictions carry $I$ affinely onto $I'$ and $J$
affinely onto $J'$.
\end{lemma}

This lemma supplies the finite set $S$ in \S\ref{sec:boundary}.

\begin{proof}
\emph{Step 1.} The complement of $I\cup J$ consists of three gaps
$[0,\min I]$, $[\max I,\min J]$, $[\max J,1]$, each of positive length
(internal, $I<J$) with dyadic endpoints; likewise for $I',J'$.

\emph{Step 2.} Cut each of the six gaps into basic intervals, e.g.\ the
cells of a sufficiently fine uniform dyadic grid.

\emph{Step 3.} Bisecting a basic interval gives two basic intervals and
raises the count by one. Bisecting on the side with fewer cells, make
the $k$-th source gap and the $k$-th target gap have equally many cells
($k=1,2,3$).

\emph{Step 4.} Now $[0,1]$ is cut on both sides into the same number of
basic intervals, with $I,J$ in the same positions as $I',J'$. Let $h$
map each source interval affinely and increasingly onto the
corresponding target interval. Adjacent pieces agree at shared
endpoints, so $h$ is an increasing homeomorphism of $[0,1]$; its
breakpoints are dyadic and each slope is a ratio of two basic lengths,
a power of two. Hence $h\in F$, and it is affine from $I$ onto $I'$ and
from $J$ onto $J'$.
\end{proof}

\emph{Covariance.} Let $h\in F$ satisfy $h(I)=I'$ with $h|_I$ affine.
Then $h\circ s_I=s_{I'}$ (both are the increasing affine bijection
$[0,1]\to I'$). If $gT^{(n)}$ and $(hg)T^{(n)}$ are basic and respect
$I$ and $I'$ respectively, then
\begin{equation}\label{dyad:covariance}
  \bigl((hg)T^{(n)}\bigr)_{I'}=(gT^{(n)})_I
\end{equation}
as sets of intervals. Indeed, the cells of $(hg)T^{(n)}$ are the images
$h(C)$ of the cells $C$ of $gT^{(n)}$; since $h$ is a bijection with
$h(I)=I'$, $h(C)\subseteq I'$ exactly when $C\subseteq I$; and
$s_{I'}^{-1}(h(C))=s_I^{-1}(C)$ for such $C$.

\subsection{The recursive coloring}\label{sec:coloring}

Fix $H$, $B$, $f$ as in Lemma~\ref{lem:displacement}, and $L>0$,
$\delta>0$ with
\begin{equation}\label{color:map}
  \norm{f(x)-f(y)}\leq L\norm{x-y},\qquad
  \norm{f(x)-x}\geq\delta\quad(x,y\in B).
\end{equation}
Choose an integer $D\ge2$ with
\begin{equation}\label{color:parameters}
  D>\frac{4L^2}{\delta^2},
\end{equation}
and internal basic intervals $I_1<\cdots<I_D$; for instance, with
$2^r>2D$, $I_j=[(2j-1)2^{-r},2j2^{-r}]$.

Define $p(T)\in B$ for every basic partition $T$ by induction on its
number of cells:
\begin{equation}\label{color:recursion}
  p(T)=
  \begin{cases}
  \displaystyle f\Bigl(\frac1D\sum_{j=1}^Dp(T_{I_j})\Bigr),
       &\text{if $T$ respects every $I_j$},\\[6pt]
  0,   &\text{otherwise}.
  \end{cases}
\end{equation}
This is well defined: in the first case $T$ has a cell outside
$I_j\ne[0,1]$, so $T_{I_j}$ has strictly fewer cells and its color is
already defined; the mean of points of the convex set $B$ lies in $B$,
the domain of $f$; and $f$ maps into $B$.

\emph{Sanity check.} For the displayed $I_j$: $p(T^{(0)})=0$ (second
case); $(T^{(r)})_{I_j}=T^{(0)}$, so $p(T^{(r)})=f(0)$;
$(T^{(2r)})_{I_j}=T^{(r)}$, so $p(T^{(2r)})=f(f(0))$.

Let
\[
  \mathcal I=\{I_i:1\leq i\leq D\}\cup\{I_i\cdot I_j:1\leq i,j\leq D\},
\]
a finite family of internal basic intervals. Call $(n,g)$
\emph{admissible} if $gT^{(n)}$ is a basic partition respecting every
member of $\mathcal I$. For all $n$, $g\in F$, $I\in\mathcal I$ put
\begin{equation}\label{color:global}
  X_I^{(n)}(g)=
  \begin{cases}
    p\bigl((gT^{(n)})_I\bigr),&\text{if $(n,g)$ is admissible},\\
    0,&\text{otherwise},
  \end{cases}
  \qquad
  \varphi_{I,J}^{(n)}(g)=\ip{X_I^{(n)}(g)}{X_J^{(n)}(g)}.
\end{equation}
These are defined on all of $F$, with $X_I^{(n)}(g)\in B$ and
$|\varphi_{I,J}^{(n)}|\le1$. By Lemma~\ref{dyad:eventual} (with
$\mathcal J=\mathcal I$), for every finite $K\subset F$ there is one $n$
with $(n,g)$ admissible for all $g\in K$.

\subsection{A uniform boundary bound}\label{sec:boundary}

For every ordered pair $I<J$ in $\mathcal I$ fix, by
Lemma~\ref{dyad:transport}, an $h_{I,J}\in F$ carrying $I,J$ affinely
onto $I_1,I_2$. Let $S$ be the finite set of these elements. It depends
only on $D$ and the $I_j$, not on any set $A$.

\emph{Finite-average comparison.} For nonempty finite $A\subset F$,
$h\in F$ and $\varphi:F\to[-1,1]$,
\begin{equation}\label{eq:finite-translation}
 \Bigl|\frac1{|A|}\sum_{g\in A}\varphi(hg)
       -\frac1{|A|}\sum_{g\in A}\varphi(g)\Bigr|
 \le\frac{|hA\mathbin{\triangle}A|}{|A|}.
\end{equation}
Indeed $\sum_{g\in A}\varphi(hg)=\sum_{x\in hA}\varphi(x)$; the terms
with $x\in hA\cap A$ cancel against the second sum, leaving
$|hA\setminus A|+|A\setminus hA|$ terms of absolute value at most one.

\begin{proposition}\label{avg:boundary}
For every nonempty finite set $A\subset F$,
\begin{equation}\label{avg:boundary-bound}
  \max_{h\in S}\frac{|hA\mathbin{\triangle}A|}{|A|}
  \geq \frac{\delta^2/L^2-4/D}{4(1-1/D)}>0.
\end{equation}
\end{proposition}

\emph{Plan.} Steps 1--2 establish properties (b) and (a) of the
overview; Step 3 is a pointwise inequality from \eqref{color:map};
Step 4 bounds an averaged variance; Step 5 combines.

\begin{proof}
\emph{Step 0 (setup).} Fix $A$ and put
\begin{equation}\label{avg:defect}
  \eta=\max_{h\in S}\frac{|hA\mathbin{\triangle}A|}{|A|}.
\end{equation}
The set $A\cup SA$, $SA=\bigcup_{h\in S}hA$, is finite, so there is one
$n$ with $(n,g)$ admissible for all $g\in A\cup SA$. This $n$ depends on
$A$ and is fixed for the rest of the proof; \emph{all} colors and
correlations below are taken at this level and the superscript is
omitted. Write
\[
  \mathbb E_A\psi=\frac1{|A|}\sum_{g\in A}\psi(g),\qquad
  \alpha=\mathbb E_A\varphi_{I_1,I_2}.
\]
The number $\alpha$ depends on $A$ and $n$; it will cancel.

\emph{Step 1 (all separated correlations are near $\alpha$).} Let $I<J$
in $\mathcal I$, $h=h_{I,J}$, $g\in A$. Both $(n,g)$ and $(n,hg)$ are
admissible, so $gT^{(n)}$ is basic and respects $I,J$, and
$(hg)T^{(n)}$ is basic and respects $I_1,I_2$. By
\eqref{dyad:covariance} (for $I\to I_1$ and for $J\to I_2$) and
\eqref{color:global},
\[
  X_{I_1}(hg)=p\bigl(((hg)T^{(n)})_{I_1}\bigr)=p\bigl((gT^{(n)})_I\bigr)=X_I(g),
  \qquad X_{I_2}(hg)=X_J(g).
\]
Hence $\varphi_{I,J}(g)=\varphi_{I_1,I_2}(hg)$ for $g\in A$. Averaging
and applying \eqref{eq:finite-translation} to the globally defined
$\varphi_{I_1,I_2}$, then \eqref{avg:defect}:
\begin{equation}\label{avg:correlation}
  \bigl|\mathbb E_A\ip{X_I}{X_J}-\alpha\bigr|\leq\eta
  \qquad\text{whenever $I<J$ or $J<I$ in $\mathcal I$}
\end{equation}
(the case $J<I$ by symmetry of the inner product). The separated pairs
used below are: (i) $I_k,I_l$ with $k\ne l$; (ii) $I_i\cdot I_j$,
$I_i\cdot I_k$ with $j\ne k$ ($s_{I_i}$ is increasing, so it preserves
order and gaps); (iii) $I_i\cdot I_j$, $I_k$ with $k\ne i$ (because
$I_i\cdot I_j\subseteq I_i$). Nothing is claimed for $I_i\cdot I_j$
against $I_i$, nor for an interval against itself.

\emph{Step 2 (recursion).} For $g\in A$ put
\[
  m(g)=\frac1D\sum_{k=1}^DX_{I_k}(g),\qquad
  z_i(g)=\frac1D\sum_{j=1}^DX_{I_i\cdot I_j}(g)\quad(1\leq i\leq D),
\]
points of $B$ by convexity. Since $gT^{(n)}$ respects $I_i$ and every
$I_i\cdot I_j$, \eqref{dyad:restriction} shows that
$(gT^{(n)})_{I_i}$ respects every $I_j$, with
$((gT^{(n)})_{I_i})_{I_j}=(gT^{(n)})_{I_i\cdot I_j}$. So the first case
of \eqref{color:recursion} applies to $T=(gT^{(n)})_{I_i}$:
\begin{equation}\label{avg:recursive-identity}
  X_{I_i}(g)=f(z_i(g))\qquad(g\in A,\ 1\leq i\leq D).
\end{equation}

\emph{Step 3 (pointwise inequality).} Fix $g\in A$ and drop it from the
notation. By \eqref{avg:recursive-identity}, $m=\frac1D\sum_if(z_i)$, so
\begin{align}
  \delta^2&\leq\norm{m-f(m)}^2
   =\Bigl\|\frac1D\sum_{i=1}^D\bigl(f(z_i)-f(m)\bigr)\Bigr\|^2\notag\\
  &\leq\frac1D\sum_{i=1}^D\norm{f(z_i)-f(m)}^2
   \leq\frac{L^2}{D}\sum_{i=1}^D\norm{z_i-m}^2,
       \label{avg:displacement}
\end{align}
using, in order, the displacement bound in \eqref{color:map}, convexity
of $\norm{\cdot}^2$, and the Lipschitz bound in \eqref{color:map}.

\emph{Step 4 (variance).} Expand
$\norm{m}^2=D^{-2}\sum_{k,l}\ip{X_{I_k}}{X_{I_l}}$. The $D(D-1)$ terms
with $k\ne l$ are of type (i), each with average at most $\alpha+\eta$
by \eqref{avg:correlation}; the $D$ diagonal terms are at most one.
The same holds for $\norm{z_i}^2$ with type (ii). So
\begin{align}
  \mathbb E_A\norm{m}^2
       &\leq\Bigl(1-\frac1D\Bigr)(\alpha+\eta)+\frac1D,
       \label{avg:parent-bound}\\
  \mathbb E_A\norm{z_i}^2
       &\leq\Bigl(1-\frac1D\Bigr)(\alpha+\eta)+\frac1D.
       \label{avg:child-bound}
\end{align}
In
$\ip{z_i}{m}=D^{-2}\sum_{j,k}\ip{X_{I_i\cdot I_j}}{X_{I_k}}$ the
$D(D-1)$ terms with $k\ne i$ are of type (iii), each with average at
least $\alpha-\eta$; the $D$ terms with $k=i$ (a child against its own
parent) are at least $-1$. So
\begin{equation}\label{avg:mixed-bound}
  \mathbb E_A\ip{z_i}{m}
    \geq\Bigl(1-\frac1D\Bigr)(\alpha-\eta)-\frac1D.
\end{equation}
These bounds hold for $\alpha,\eta$ of any sign or size. From
$\norm{z_i-m}^2=\norm{z_i}^2+\norm{m}^2-2\ip{z_i}{m}$,
\begin{equation}\label{avg:variance}
  \mathbb E_A\norm{z_i-m}^2
   \le2\Bigl(1-\frac1D\Bigr)(\alpha+\eta)+\frac2D
      -2\Bigl(1-\frac1D\Bigr)(\alpha-\eta)+\frac2D
   =\frac4D+4\Bigl(1-\frac1D\Bigr)\eta .
\end{equation}

\emph{Step 5 (conclusion).} Average \eqref{avg:displacement} over
$g\in A$ and insert \eqref{avg:variance}:
\[
  \delta^2\leq\frac{L^2}{D}\sum_{i=1}^D\mathbb E_A\norm{z_i-m}^2
  \leq L^2\Bigl(\frac4D+4\Bigl(1-\frac1D\Bigr)\eta\Bigr).
\]
Since $D\ge2$, solving for $\eta$ gives \eqref{avg:boundary-bound}; the
right-hand side is positive by \eqref{color:parameters} and involves
neither $A$ nor $n$.
\end{proof}

\begin{proof}[Proof of Theorem~\ref{thm:main}]
Let $c_0>0$ be the right-hand side of \eqref{avg:boundary-bound}. If
$F$ were amenable, \eqref{eq:folner} with the finite set $S$ and
$\varepsilon=c_0$ would give a nonempty finite $A$ with
$\max_{h\in S}|hA\triangle A|/|A|<c_0$, contradicting
Proposition~\ref{avg:boundary}.
\end{proof}

\begin{remark}[What is used about $F$]
The proof of Proposition~\ref{avg:boundary} uses $F$ only through
Step 1 (transport identities on $A$) and Step 2 (recursion on $A$),
at one level $n$. Sanity check: if every $h_{I,J}$ were the identity,
Step 1 would force all $X_I$ to coincide, and Step 2 would give a
fixed point of $f$, which is impossible; this matches
\eqref{avg:boundary-bound}, which would then read $0\ge c_0$. The
content specific to $F$ is that both properties hold with the
transports of Lemma~\ref{dyad:transport}, which change the size and
position of intervals, while the recursion stays well founded.
\end{remark}

\section{Consequences}\label{sec:consequences}

Two corollaries are deduced in the source from companion results;
neither is used above.

\begin{corollary}\label{cor:nonunitarizable}
For every \(\varepsilon>0\), there are a separable complex Hilbert space
\(\mathcal K\) and a representation \(\pi:F\to\mathrm{GL}(\mathcal K)\)
such that \(\sup_{g\in F}\norm{\pi(g)}\le1+\varepsilon\), but no bounded
invertible operator \(V\) on \(\mathcal K\) makes every
\(V\pi(g)V^{-1}\) unitary.
\end{corollary}

\begin{proof}
$F$ is countable (each element is given by finitely many dyadic
breakpoints and slopes) and nonamenable by Theorem~\ref{thm:main}; the
source concludes by the companion unitarizability theorem
\cite[Theorem~1.1]{OpenAIUnitarizability2026}.
\unverified{The source does not reproduce the statement of the companion
theorem, so I could not check that countability and nonamenability are
its only hypotheses or that it yields the bound $1+\varepsilon$. A
theorem of that generality would settle Dixmier's unitarizability
problem, which I know only as open.}
\end{proof}

For a finite symmetric generating set
$\Sigma\subset F\setminus\{\id\}$, let $\mathcal G_\Sigma$ be the simple
undirected Cayley graph with edges $\{x,xs\}$, $s\in\Sigma$. In
Bernoulli bond percolation each edge is independently open with
probability $p$. Let $p_c$ be the threshold for existence of an infinite
open cluster and $p_u$ the threshold for almost-sure uniqueness of the
infinite open cluster. The connection kernel
$\mathcal T_p(x,y)=\mathbb P_p(x\leftrightarrow y)$ acts on
counting-measure $\ell^2(F)$, initially on finitely supported functions
by summation over $y$; put
$p_{2\to2}=\sup\{p\in[0,1]:\mathcal T_p\text{ is bounded on }\ell^2(F)\}$.

\begin{corollary}\label{cor:percolation}
For every such generating set \(\Sigma\),
\[
  \norm{\mathcal T_{p_c}}_{2\to2}<\infty,\qquad p_c<p_{2\to2}\le p_u.
\]
For each fixed \(p\in(p_c,p_u)\), there are almost surely infinitely many
infinite open clusters. Moreover, there are deterministic
\(p_c<p_1<p_2<1\), depending on \(\Sigma\), such that in the coupling
with independent uniform \([0,1]\) edge labels \(U_e\), opening \(e\)
when \(U_e\le p\), there are almost surely infinitely many infinite open
clusters simultaneously for every \(p\in[p_1,p_2]\).
\end{corollary}

\begin{proof}
The source applies the companion percolation result
\cite[Corollary~1.2]{OpenAIPercolation2026} to $F$, nonamenable by
Theorem~\ref{thm:main}.
\unverified{The statement of the companion corollary is not reproduced
in the source, so none of the conclusions can be checked from this
document. The general inequality $p_c<p_u$ for nonamenable Cayley graphs
is, to my knowledge, an open conjecture.}
\end{proof}

\appendix
\section{An explicit map with positive displacement}\label{app:auxiliary}

This appendix proves the Hilbert-space case of
Lemma~\ref{lem:displacement} \cite{BenyaminiSternfeld1983} directly. It
is independent of \S\ref{sec:proof}; here $r$ denotes a distance.

\begin{proposition}\label{aux:map}
Let $H=L^2([0,1];\R^2)$ be a real Hilbert space and let
$B=\{x\in H:\norm{x}\le1\}$. There is a globally Lipschitz map $f:B\to B$
such that
\[
  \norm{f(x)}=1,\qquad\norm{f(x)-x}\ge\frac12\quad(x\in B).
\]
\end{proposition}

\emph{Plan.} We build a Lipschitz $G:B\to H$ with
$\norm{G(x)}\ge\rho/4$ and $G(x)=x$ for $\norm{x}\ge1/2$, and set
$f=-G/\norm{G}$. $G$ is the identity outside a thin tube around a curve
$\gamma$ through the origin. The curve enters along a straight ray and
then spirals forever inside a small ball without approaching itself
(Lemma~\ref{aux:tube}); inside the tube, $G$ moves points along the
curve direction so that $0$ is avoided (Lemma~\ref{aux:rotation}
supplies a needed rotation).

Define
\[
 a(t)=
 \begin{cases}
 t/8,&t\le2,\\
 5/16-(3-t)^2/16,&2<t<3,\\
 5/16,&t\ge3,
 \end{cases}
 \qquad
 \theta(t)=
 \begin{cases}
 0,&t\le1,\\
 (t-1)^2/2,&1<t<2,\\
 t-3/2,&t\ge2,
 \end{cases}
\]
and, for $w\in[0,1]$,
\[
 v(t)(w)=(\cos(\theta(t)w),\sin(\theta(t)w)),\qquad\gamma(t)=a(t)v(t).
\]

\emph{Elementary facts.} The pieces match to first order at the
junctions ($a(2)=1/4$, $a'(2)=1/8$, $a'(3)=0$; $\theta(2)=1/2$,
$\theta'(1)=0$, $\theta'(2)=1$), so $a,\theta$ are $C^1$ with
Lipschitz derivatives, $0\le a'\le1/8$, $0\le\theta'\le1$. Moreover
$a\le5/16$; $a(t)\ge1/8$ for $t\ge1$ and $a(t)\ge1/4$ for $t\ge2$;
$\theta=0$ on $(-\infty,1]$ and $\theta$ is strictly increasing on
$[1,\infty)$. Also $\norm{v(t)}=1$, $\norm{\gamma(t)}=|a(t)|$,
$\gamma(0)=0$, and for $t\le1$: $v(t)=(1,0)$ and
$\gamma(t)=(t/8)(1,0)$.

\begin{lemma}\label{aux:tube}
The curve $\gamma$ is injective, has a bounded Lipschitz derivative, and
\[
 c:=\inf_{t\in\R}\norm{\gamma'(t)}=\frac18.
\]
Its unit tangent $u(t)=\gamma'(t)/\norm{\gamma'(t)}$ is Lipschitz,
$\ip{u(t)}{v(t)}\ge0$, and $u(t)=v(t)$ for $t\le1$.
There is a radius $0<\rho<1/32$ such that every point in
\[
 U=\{x\in H:\dist(x,\gamma(\R))<\rho\}
\]
has a unique nearest point $\gamma(t(x))$. Writing
$e(x)=x-\gamma(t(x))$, we have $e(x)\perp u(t(x))$, and
\begin{equation}\label{aux:parameter}
 |t(x)-t(y)|\le16\norm{x-y}
 \quad\text{if }x,y\in U\text{ and }\norm{x-y}<\rho.
\end{equation}
\end{lemma}

\begin{proof}
\emph{Step 1 (derivative and speed).} Let
$v^\perp(t)(w)=(-\sin(\theta(t)w),\cos(\theta(t)w))$. For each $w$,
\[
 \partial_t\bigl(\gamma(t)(w)\bigr)
 =a'(t)v(t)(w)+a(t)\theta'(t)\,w\,v^\perp(t)(w).
\]
This is locally Lipschitz in $t$ uniformly in $w\in[0,1]$, so the
difference quotients converge uniformly in $w$, hence in $H$; the right
side is $\gamma'(t)$. As $v\perp v^\perp$ pointwise and
$\int_0^1w^2dw=1/3$,
\[
 \norm{\gamma'(t)}^2=a'(t)^2+\frac{a(t)^2\theta'(t)^2}{3},
 \qquad\ip{\gamma'(t)}{v(t)}=a'(t).
\]
For $t\le2$, $a'=1/8$, so $\norm{\gamma'}\ge1/8$, with equality for
$t\le1$ (where $\theta'=0$). For $t\ge2$, $a\ge1/4$ and $\theta'=1$, so
$\norm{\gamma'}^2\ge1/48>1/64$. Hence $c=1/8$.

\emph{Step 2 (regularity).} Pointwise
$|v(t)(w)-v(s)(w)|\le|\theta(t)-\theta(s)|w$, so
$\norm{v(t)-v(s)}\le|t-s|$; likewise for $wv^\perp$. The coefficients
$a'$ and $a\theta'$ are bounded and Lipschitz ($a\theta'$ vanishes for
$t\le1$, and for $t\ge1$ both factors are bounded and Lipschitz). So
$\gamma'$ is bounded, by $M$ say, and Lipschitz, with constant
$K_\gamma$. Since $\norm{\gamma'}\ge1/8$, $u$ is Lipschitz, with
constant $K_u$. Finally $\ip uv=a'/\norm{\gamma'}\ge0$, and for $t\le1$,
$\gamma'=v/8$, so $u=v$.

\emph{Step 3 (injectivity).} With $\sinc(q)=\sin(q)/q$, $\sinc(0)=1$,
\[
 \ip{v(s)}{v(t)}=\int_0^1\cos\bigl((\theta(t)-\theta(s))w\bigr)dw
 =\sinc(\theta(t)-\theta(s)).
\]
Since $|\sinc(q)|<1$ for $q\ne0$, the unit vectors $v(s),v(t)$ are
linearly independent whenever $\theta(s)\ne\theta(t)$ (equality case of
Cauchy--Schwarz). Let $s<t$. If $t\le1$,
$\gamma(s)=(s/8)(1,0)\ne\gamma(t)$. If $t>1$, then
$\theta(s)<\theta(t)$, so $a(s)v(s)=a(t)v(t)$ would force $a(t)=0$,
contradicting $a(t)\ge1/8$.

\emph{Step 4 (uniform separation).} We claim
\begin{equation}\label{aux:separation}
 \sigma(d):=\inf_{|t-s|\ge d}\norm{\gamma(t)-\gamma(s)}>0\qquad(d>0).
\end{equation}
Otherwise there are $s_n<t_n$ with $t_n-s_n\ge d$ and
$\norm{\gamma(t_n)-\gamma(s_n)}\to0$; pass to a subsequence with
$s_n\to s_*$, $t_n\to t_*$ in $[-\infty,\infty]$. Each case is
contradictory:
\begin{itemize}
\item Both finite: $t_*-s_*\ge d$ and $\gamma(s_*)=\gamma(t_*)$,
against Step 3.
\item $t_*=-\infty$: eventually $s_n<t_n\le1$ and the distance is
$(t_n-s_n)/8\ge d/8$.
\item $s_*=-\infty<t_*$: $\norm{\gamma(s_n)}=|s_n|/8\to\infty$ while
$\norm{\gamma(t_n)}$ stays bounded.
\item $s_*$ finite, $t_*=+\infty$: $\theta(t_n)-\theta(s_n)\to\infty$,
so $\ip{v(s_n)}{v(t_n)}\to0$ and
$\norm{\gamma(t_n)-\gamma(s_n)}^2\to(5/16)^2+a(s_*)^2>0$.
\item $s_*=+\infty$: eventually $s_n,t_n\ge3$, so $a=5/16$,
$\theta(t_n)-\theta(s_n)=t_n-s_n\ge d$, and
\[
 \norm{\gamma(t_n)-\gamma(s_n)}^2
 =2(5/16)^2\bigl(1-\sinc(t_n-s_n)\bigr)
 \ge2(5/16)^2\bigl(1-\sup_{q\ge d}\sinc(q)\bigr)>0,
\]
since $\sinc$ is continuous, less than $1$ on $[d,\infty)$, and tends
to $0$.
\end{itemize}

\emph{Step 5 (tangent estimate; choice of $\rho$).} Fix $d>0$ with
$K_\gamma d\le c/2$. Integrating $\gamma'(\tau)-\gamma'(s)$ from $s$ to
$t$,
\[
 \norm{\gamma(t)-\gamma(s)-(t-s)\gamma'(s)}\le\frac{K_\gamma}2|t-s|^2 .
\]
Pair with $u(s)$ and use $\ip{\gamma'(s)}{u(s)}=\norm{\gamma'(s)}\ge c$:
\begin{equation}\label{aux:tangent}
 |\ip{\gamma(t)-\gamma(s)}{u(s)}|
 \ge c|t-s|-\frac{K_\gamma}{2}|t-s|^2\ge\frac{3c}{4}|t-s|
 \qquad(|t-s|\le d).
\end{equation}
Choose $\rho>0$ with
\begin{equation}\label{aux:radius}
 \rho<\frac1{32},\qquad4\rho<\sigma(d),\qquad\rho K_u<\frac c4.
\end{equation}

\emph{Step 6 (a nearest point exists; orthogonality).} Let $x\in U$ and
let $\gamma(t_n)$ be a minimizing sequence for
$\dist(x,\gamma(\R))$. Late terms satisfy
$\norm{x-\gamma(t_n)}<2\rho$, so they are pairwise less than
$4\rho<\sigma(d)$ apart, and by \eqref{aux:separation} their parameters
differ by less than $d$. So the late $t_n$ lie in a compact interval,
and a convergent subsequence attains the minimum. At a minimizing $t$,
the derivative of $\norm{x-\gamma(t)}^2$ vanishes:
$(x-\gamma(t))\perp\gamma'(t)$, i.e.\ $\perp u(t)$.

\emph{Step 7 (parameter estimate; uniqueness).} Let $x,y\in U$,
$\norm{x-y}<\rho$, with minimizing parameters $t,s$ and
$x=\gamma(t)+e_t$, $y=\gamma(s)+e_s$, $\norm{e_t},\norm{e_s}<\rho$.
Then $\norm{\gamma(t)-\gamma(s)}<3\rho<\sigma(d)$, so $|t-s|<d$. Since
$e_s\perp u(s)$ and $e_t\perp u(t)$,
\[
 \ip{x-y}{u(s)}=\ip{\gamma(t)-\gamma(s)}{u(s)}+\ip{e_t}{u(s)-u(t)},
\]
so by \eqref{aux:tangent} and \eqref{aux:radius}
\[
 \norm{x-y}\ge\Bigl(\frac{3c}{4}-\rho K_u\Bigr)|t-s|
 \ge\frac c2|t-s|=\frac1{16}|t-s|.
\]
Taking $y=x$ gives uniqueness of the minimizing parameter; in general
this is \eqref{aux:parameter}.
\end{proof}

\begin{lemma}\label{aux:rotation}
There is a family of orthogonal operators $R_t:H\to H$, Lipschitz
in operator norm, such that $R_tu(t)=v(t)$ and $R_t$ is the identity
for $t\le1$.
\end{lemma}

It is used to turn a residual $e\perp u(t)$ into a vector
$R_te\perp v(t)$.

\begin{proof}
Suppress $t$ and put $\kappa=\ip uv\ge0$ (Lemma~\ref{aux:tube}). Define
\[
 Ah=v\ip uh-u\ip vh,\qquad R=\mathrm{Id}+A+\frac{A^2}{1+\kappa}.
\]
$A$ is skew-adjoint, since
$\ip{Ah}{k}=\ip vk\ip uh-\ip uk\ip vh=-\ip h{Ak}$. If $u,v$ are
independent, compute on their span: $Au=v-\kappa u$,
$Av=\kappa v-u$, hence $A^2u=-(1-\kappa^2)u$, $A^2v=-(1-\kappa^2)v$, so
$R=\kappa\mathrm{Id}+A$ there. Then
$Ru=v$, and
$(\kappa-A)(\kappa+A)=\kappa^2-A^2=\mathrm{Id}$ on the span, so $R$ is
orthogonal there; on the orthogonal complement $A=0$ and
$R=\mathrm{Id}$. If $u,v$ are dependent then $u=v$ (as $\kappa\ge0$),
$A=0$, $R=\mathrm{Id}$; this covers $t\le1$. Finally $A_t$ and
$\kappa(t)$ are Lipschitz because $u,v$ are, $\norm{A_t}\le2$, and
$1+\kappa\ge1$, so $t\mapsto R_t$ is Lipschitz in operator norm, with
constant $K_R$ say.
\end{proof}

\begin{proof}[Proof of Proposition~\ref{aux:map}]
Fix $\rho$, $U$, $t(\cdot)$, $e(\cdot)$ from Lemma~\ref{aux:tube} and
$R_t$ from Lemma~\ref{aux:rotation}.

\emph{Step 1 (definition of $G$).} Set
\[
 b(t)=\min\{a(t),-1/8\},\qquad
 q(t)=b(t)-a(t)=-\max\{a(t)+1/8,0\},
\]
so $|q|\le5/16+1/8=7/16$, $\Lip(q)\le1/8$, and $q(t)=0$ when
$a(t)\le-1/8$. Let $\beta,\chi:[0,\infty)\to[0,1]$ be
\[
 \beta(r)=
 \begin{cases}
 1,&r\le\rho/4,\\
 2-4r/\rho,&\rho/4<r<\rho/2,\\
 0,&r\ge\rho/2,
 \end{cases}
 \qquad
 \chi(r)=
 \begin{cases}
 1,&r\le\rho/2,\\
 2-2r/\rho,&\rho/2<r<\rho,\\
 0,&r\ge\rho.
 \end{cases}
\]
Let $G(x)=x$ for $x\in B\setminus U$. For $x\in B\cap U$ write
$t=t(x)$, $e=e(x)$, $r=\norm e=\dist(x,\gamma(\R))$, and set
\begin{equation}\label{aux:G}
 \begin{split}
 G(x)&=\gamma(t)+\beta(r)q(t)v(t)
       +\bigl((1-\chi(r))\mathrm{Id}+\chi(r)R_t\bigr)e\\
 &=x+\beta(r)q(t)v(t)+\chi(r)(R_t-\mathrm{Id})e.
 \end{split}
\end{equation}

\emph{Step 2 ($G$ is Lipschitz).} Let $x,y\in B$, $\Delta=\norm{x-y}$.

\emph{(i) $x,y\in U$, $\Delta<\rho$.} By \eqref{aux:parameter},
$|t(x)-t(y)|\le16\Delta$; hence
$\norm{e(x)-e(y)}\le\Delta+\norm{\gamma(t(x))-\gamma(t(y))}
\le(1+16M)\Delta$; and $r$ is $1$-Lipschitz. The factors in
\eqref{aux:G} therefore satisfy:
\begin{center}
\begin{tabular}{lcccccc}
factor & $\beta(r)$ & $q(t)$ & $v(t)$ & $\chi(r)$ & $R_t-\mathrm{Id}$ & $e$\\
bound & $1$ & $7/16$ & $1$ & $1$ & $2$ & $\rho$\\
change $\le$ & $4\Delta/\rho$ & $2\Delta$ & $16\Delta$ & $2\Delta/\rho$ & $16K_R\Delta$ & $(1+16M)\Delta$
\end{tabular}
\end{center}
A product of three factors changes by at most the sum, over the
factors, of that factor's change times the bounds of the other two.
Hence
\[
 \norm{G(x)-G(y)}\le\Bigl[1+\Bigl(\frac7{4\rho}+2+7\Bigr)
   +\bigl(4+16K_R\rho+2(1+16M)\bigr)\Bigr]\Delta .
\]

\emph{(ii) $\Delta\ge\rho$.} On $B$,
$\norm{G(x)}\le1+7/16+2\rho<3/2$, so
$\norm{G(x)-G(y)}\le3\le(3/\rho)\Delta$.

\emph{(iii) $x\in U$, $y\notin U$, $\Delta<\rho$.} Both cutoffs satisfy
$\beta(r)\le\chi(r)\le2(\rho-r)/\rho$, so
\[
 \norm{G(x)-x}\le\frac{2(\rho-r)}\rho\Bigl(\frac7{16}+2\rho\Bigr)
 =\Bigl(\frac{7}{8\rho}+4\Bigr)(\rho-r).
\]
As $r$ is $1$-Lipschitz and $r(y)\ge\rho$, $\rho-r\le\Delta$. With
$G(y)=y$: $\norm{G(x)-G(y)}\le(7/(8\rho)+5)\Delta$.

\emph{(iv) $x,y\notin U$.} $G$ is the identity.

\emph{Step 3 (nonvanishing).} We claim
\begin{equation}\label{aux:nonzero}
 \norm{G(x)}\ge\rho/4\qquad(x\in B).
\end{equation}
Outside $U$: $\norm{G(x)}=\norm{x-\gamma(0)}\ge\rho$. Inside $U$:
\begin{itemize}
\item $r\le\rho/2$. Then $\chi=1$ and
$G(x)=(a(t)+\beta(r)q(t))v(t)+R_te$, with
$\ip{R_te}{v(t)}=\ip{R_te}{R_tu(t)}=\ip e{u(t)}=0$ and
$\norm{R_te}=r$. If $r\le\rho/4$, then $\beta=1$ and the coefficient of
$v(t)$ is $b(t)\le-1/8$, so $\norm{G(x)}\ge1/8>\rho/4$. If
$r\ge\rho/4$, then $\norm{G(x)}\ge\norm{R_te}\ge\rho/4$.
\item $\rho/2\le r<\rho$. Then $\beta=0$ and the last term in the first
line of \eqref{aux:G} has norm at most $r<\rho$. If $t>1$,
$\norm{\gamma(t)}=a(t)\ge1/8$, so
$\norm{G(x)}\ge1/8-\rho>\rho/4$. If $t\le1$, $R_t=\mathrm{Id}$, so
$G(x)=a(t)v(t)+e$ with $e\perp u(t)=v(t)$, and
$\norm{G(x)}\ge r\ge\rho/2$.
\end{itemize}

\emph{Step 4 ($G(x)=x$ when $\norm x\ge1/2$).} Outside $U$ this is the
definition. Inside $U$,
\[
 |a(t)|=\norm{\gamma(t)}\ge\norm x-r>1/2-\rho>15/32>5/16 .
\]
Since $a\le5/16$ everywhere, $a(t)<-5/16$; thus $a(t)\le-1/8$ and
$t\le1$, so $q(t)=0$ and $R_t=\mathrm{Id}$, and \eqref{aux:G} gives
$G(x)=x$.

\emph{Step 5 (the map $f$).} Define $f(x)=-G(x)/\norm{G(x)}$. The map
$y\mapsto y/\norm y$ is Lipschitz with constant at most $8/\rho$ on
$\{\norm y\ge\rho/4\}$, so $f$ is Lipschitz by Step 2 and
\eqref{aux:nonzero}, and $\norm{f(x)}=1$. If $\norm x<1/2$, then
$\norm{f(x)-x}\ge1-\norm x>1/2$. Otherwise $G(x)=x$ by Step 4,
$f(x)=-x/\norm x$, and $\norm{f(x)-x}=1+\norm x\ge3/2$.
\end{proof}

\section*{Verification checklist}

\begin{enumerate}
\item \textbf{Order of choices.} $f,L,\delta,D$, the intervals $I_j$ and
the set $S$ are fixed before $A$; only $n$ and $\alpha$ depend on $A$,
and one $n$ is used for every color in the proof
(Proposition~\ref{avg:boundary}, Step 0).
\item \textbf{Transport identity.} $X_I(g)=X_{I_1}(h_{I,J}g)$ needs
admissibility of both $g$ and $h_{I,J}g$ and the exact covariance
\eqref{dyad:covariance} (Step 1; \S\ref{sec:dyadic}).
\item \textbf{Which pairs are separated.} Types (i)--(iii) in Step 1,
and the counts $D(D-1)$ versus $D$ in Step 4; a child against its own
parent is only bounded by $-1$.
\item \textbf{Recursion identity} \eqref{avg:recursive-identity}:
admissibility plus \eqref{dyad:restriction} put
$(gT^{(n)})_{I_i}$ in the first case of \eqref{color:recursion}
(Step 2); well-foundedness of \eqref{color:recursion}
(\S\ref{sec:coloring}).
\item \textbf{Cancellation of $\alpha$} in \eqref{avg:variance} and the
arithmetic from Step 5 to \eqref{avg:boundary-bound}, using
\eqref{color:parameters}.
\item \textbf{Transports lie in $F$} (Lemma~\ref{dyad:transport},
Step 4) and fine images are basic (Lemma~\ref{dyad:eventual}, Step 2).
\item \textbf{Appendix.} Uniform separation \eqref{aux:separation},
especially the case of both parameters tending to $+\infty$
(Lemma~\ref{aux:tube}, Step 4); nonvanishing \eqref{aux:nonzero}
(Proposition~\ref{aux:map}, Step 3); the Lipschitz case analysis
(Step 2).
\item \textbf{Corollaries.} They rest on companion results not
reproduced in the source; see the boxed notes in
\S\ref{sec:consequences}.
\end{enumerate}

\bibliographystyle{plain}
\bibliography{sources}
\end{document}
